\section{Resumen}
\subsection{Propósito y alcance del estudio}

Este estudio compara experimentalmente cuatro métodos de búsqueda de patrones
en texto: búsqueda por fuerza bruta, Suffix Array,
Knuth--Morris--Pratt (KMP) y Rabin--Karp. La fuerza bruta se incluye como
\textbf{línea base} para contextualizar el desempeño de los métodos estudiados.
Para efectos de los experimentos, \textit{búsqueda} significa encontrar y
contar todas las ocurrencias del patrón en el texto, no devolver sus posiciones.
Se analiza cómo influyen el largo del texto, el largo del patrón y el número de
ocurrencias en el tiempo de búsqueda; además, se consideran el costo de
construcción del índice y el consumo de memoria. Los resultados se interpretan
junto con las complejidades teóricas para identificar los escenarios en que
conviene cada alternativa.

\subsection{Estructuras y algoritmos estudiados}

\textbf{Fuerza bruta:} Revisa cada posición posible del texto y compara desde
allí los caracteres con el patrón. No requiere preprocesamiento y sirve como
referencia sencilla, aunque puede repetir comparaciones.

\textbf{Suffix Array:} Ordena las posiciones iniciales de los sufijos del
texto. Para buscar un patrón, localiza mediante búsqueda binaria el intervalo
de sufijos que comienza con él y cuenta sus ocurrencias. Su índice requiere
una etapa de construcción previa.

\textbf{KMP:} Busca el patrón recorriendo el texto y usando información de prefijos del
patrón para evitar comparaciones repetidas. Su búsqueda tiene tiempo lineal en
el largo del texto y del patrón, sin requerir un índice persistente.

\textbf{Rabin--Karp:} Compara primero valores hash del patrón y de ventanas del texto;
cuando coinciden, verifica los caracteres para confirmar la ocurrencia. El
hash de las ventanas se actualiza incrementalmente y su desempeño depende,
entre otros factores, de las colisiones.
